\section{Evidence Reuse Focuses Each Iteration}

\subsection{QoR-RAG Retrieves Compatible Actions}

QoR-RAG builds a query from source metadata, baseline QoR, synthesis
diagnostics, resource headroom, rejection history, and the task description.
Its deterministic ranker weights token overlap by $3,2,4,2,3,1$ in that order.
Keyword indicators map the query to reduction, memory-port, dataflow, stream,
and factorization families. Each query returns at most one curated rule, one
verified success, and one measured failure. Measured cases must match the
target part and Vitis version. Workload guards compare the task description
with family tags such as AES, GEMM, stencil, and dot product. Provenance checks
reject paths that contain hidden, reference, or evaluator components. A
1,800-token cap bounds the retrieved context.

\subsection{Failure Reflection Converts Errors into Constraints}

Failure Reflection extracts the evidence that the next iteration needs. It
removes ANSI control text, redacts paths and secrets, deduplicates lines, and
selects bounded diagnostics that contain compiler, mismatch, timing, stream,
deadlock, or timeout signals. The resulting record contains the failure class,
candidate diff, implicated pragmas, loops, arrays, and the next constraint. A
normalized signature counts repeated failure patterns. The next prompt and
QoR-RAG query both consume this record, so Failure Reflection uses zero extra
LLM requests.

A candidate fingerprint removes comments, normalizes whitespace, and removes
optional braces around a single-statement \texttt{for} loop. The controller
uses this fingerprint to reject unchanged or repeated candidates before tool
execution. It stops on budget exhaustion, duplicate candidates, or search
stagnation.
